\documentclass[10pt,a4paper]{article}
\usepackage[utf8]{inputenc}
\usepackage[spanish]{babel}
\usepackage{amsmath}
\usepackage{amsfonts}
\usepackage{amssymb}
\usepackage{graphicx}
\usepackage{natbib}
\usepackage{lineno}
\usepackage{ragged2e}
\usepackage{multicol}
\setlength\columnsep{37pt}
\usepackage{enumerate} 
\usepackage[left=1.94cm,top=1.59cm,right=1.9cm,bottom=0.59cm]{geometry} 
\usepackage{fancyhdr}
\usepackage{url}
\usepackage{cite} 


\begin{document}
		
		\begin{center}
			\huge \textbf{Business Model Canvas aplicado a una empresa} 
		\end{center}
		\vspace{\baselineskip}
		\begin{center}
			\includegraphics[scale=0.37]{./Imagenes/logo}
		\end{center}
		\begin{multicols}{2}
			\small
			\begin{center}
				Nelia Escalante Marón\\
				2014049551\\
				UPT - Ingenierí­a de Sistemas\\
				EPIS\\
				Tacna, Perú\\
				\vspace{\baselineskip}
				Yerson Coaquira Calizaya\\
				2015053225\\
				UPT - Ingenierí­a de Sistemas\\  
				EPIS\\
				Tacna, Perú\\                 
				\vspace{\baselineskip}
				Flor Condori Gutierrez\\
				2015053227\\
				UPT - Ingenierí­a de Sistemas\\ 
				EPIS\\	
				Tacna, Perú\\                 
				\columnbreak
				
				\vspace{\baselineskip}
				Christian Cespedes Medina\\
				2010036256\\
				UUPT - Ingenierí­a de Sistemas\\  
				EPIS\\	
				Tacna, Perú\\                

				\vspace{\baselineskip}
				Javier Octavio Arteaga Ramos \\
				2007028981\\
				UPT - Ingenierí­a de Sistemas\\  
				EPIS\\	
				Tacna, Perú\\                 

			\end{center}
			\normalsize			
		\end{multicols}
		\vspace{\baselineskip}
		\begin{multicols}{2}
		\textbf{\textit{\large Abstract}}\rule[1.5mm]{5mm}{0.1mm} El modelo Canvas cuenta con 9 bloques los cuales hacen referencia a las características de la empresa que se quiere crear. Debemos tener en cuenta que al inicio puede costarnos un poco insertar los datos necesarios en cada bloque, y eso puede deberse a que el modelo de negocio aun no está bien definido.
		
		Los 9 bloques que va a llenar segun este modelo son: Propuesta de Valor, Los canales, Relaciones con los Clientes, Fuentes de ingresos, Recurso clave, Activida des principales, Alianza Clave y Estructura de Precio.
				
		\textit{The Canvas model has 9 blocks which refer to the characteristics of the company you want to create. We must take into account that at the beginning it can cost us a bit to insert the necessary data in each block, and that may be because the business model is not yet well defined.}
		
		\textit{In this final work we will apply the blocks of the BMC model to a company, these are: value proposition, channels, customer relations, revenue sources, key resource, core activities, key alliance and price structure}
		
		\vspace{\baselineskip}
			
		\textbf{\textit{\large Keybwords}}\rule[1.5mm]{5mm}{0.1mm} BMC, Web 2.0, aplicacion del Modelo Canvas a una empresa, cloud
		
		\columnbreak
		
		\section{Marco Teórico} 
		
		Todo emprendedor posee una idea de negocio, sin embargo, su puesta en práctica no resulta fácil, y mucho menos sacarle rentabilidad. Por ello contamos con varias estrategias, entre ellas el modelo Canvas. Estas estrategias buscan asegurarnos de que nuestras iniciativas llegaran a tener a éxito. Sin embargo, no todos los modelos de negocio nos dan soluciones perfectas. Es por ello por lo que surge el modelo Canvas.\\
		
		Modelo que se ha convertido desde el 2008 como herramienta estrella en la gestión estratégica y empresarial de un negocio. El modelo Canvas permite ver y moldear en un solo folio, estructurado en nueve elementos, cual es el modelo de nuestro negocio.  Y lo mejor de todo, es tan sencillo que puede ser aplicado en cualquier escenario, ya sea una pequeña, mediana y gran empresa.  Además, no sólo sirve para las nuevas empresas sino también para aquellas que ya están establecidas. \\
		
		En definitiva, este modelo trata de aprender muy rápido sobre el mercado, en un corto tiempo y con el mínimo coste.  Con el objetivo de lograr un modelo que busque la agilidad y la reducción del tiempo en el desarrollo de iniciativas empresariales, para finalmente generar productos y servicios que cumplan con las necesidades de los clientes y aporten valor.\\
		
		%Escribir por aqui(borrar esto)
		
			
			
			
		\end{multicols}	
			
\end{document}}